\subsection{Backpropagation}

Backpropagation computes these gradients efficiently by applying the chain rule from the output layer back through the network \cite{Rumelhart1986}. In compact form, it provides the quantities
$\nabla_{\theta}\mathcal{L}$
where $\theta$ denotes the trainable parameters and $\mathcal{L}$ is the loss function.

In the case of this pruning-oriented report we do not require the full derivation above. It is sufficient that back-propagation relates model error to specific parameters and structures, which will be exploited by the gradient-based pruning criteria, Taylor approximation, Fisher scores and post-pruning recovery process detailed in later chapters.

\section{Conclusion}

In this chapter, we have established the mathematical and optimization framework of deep neural networks. We described the building blocks-the affine and non-linear transformation of feed forward networks and formulated training of a feed forward network as an empirical risk minimization problem. We also described the central optimization concepts of gradient and Hessian matrices, crucial to the learning process of the networks and to pruning algorithms in assessing parameters relevance.


Though the above theory applies to any feed-forward network, the architecture used in modern NLP are often extremely specialized. To know how pruning is practically applied one must be aware of the mapping of generic matrix operations to the structures within an architecture. Therefore the subsequent chapter deals solely with the Transformer architecture and its components, how parameters are structured and the drastic bottleneck which necessitates compression.